\chapter{Dobre praktyki pracy z MoonChunk}

\section{Po co w ogole mowic o praktykach}
Sama znajomosc skladni nie wystarcza, aby pisac dobry kod. Mozna stworzyc program, ktory
technicznie dziala, ale jest trudny do czytania, zly do rozbudowy i nieprzyjemny w
utrzymaniu. Dobre praktyki sprawiaja, ze projekt pozostaje czytelny po tygodniu,
miesiacu i semestrze.

\section{Oddzielaj dane od prezentacji}
Jeden z najlepszych nawykow w MoonChunk polega na tym, aby najpierw przygotowac dane,
a dopiero potem generowac z nich HTML.

Lepszy styl:

\begin{lstlisting}[style=moonchunk,caption={Najpierw dane, potem content}]
let heading: string = "Aktualnosci";
let intro: string = "Najnowsze wpisy z projektu";

content {
  <h1>{heading}</h1>
  <p>{intro}</p>
};
\end{lstlisting}

Gorszy styl to upychanie wszystkiego bezposrednio do jednego, ogromnego bloku \mc{content}.

\section{Nazywaj rzeczy tak, jak je rozumiesz}
Nazwy typu \mc{x}, \mc{tmp}, \mc{data2}, \mc{something} szybko zamieniaja projekt w
labirynt. Dobra nazwa bardzo pomaga, nawet jesli jej wpisanie zajmie trzy sekundy dluzej.

Dobre przykłady:

\begin{itemize}
  \item \mc{siteTitle},
  \item \mc{isPublished},
  \item \mc{cardsHtml},
  \item \mc{featuredPosts},
  \item \mc{buildAuthorLabel}.
\end{itemize}

\section{Trzymaj metadata razem}
Jesli kilka stron korzysta z tych samych pol metadata, wydziel osobny chunk z konfiguracja
wspolna i dolacz go przez \mc{@include}. To jeden z najlatwiejszych sposobow na unikniecie
kopiowania kodu.

\section{Uzywaj funkcji do logiki, nie do wszystkiego}
Funkcje sa swietne, ale nie musza zastapic calej struktury projektu. Najlepiej sprawdzaja
sie wtedy, gdy:

\begin{itemize}
  \item masz powtarzalne obliczenie,
  \item chcesz nazwac fragment logiki,
  \item ten sam schemat pojawia sie w wielu miejscach,
  \item chcesz odchudzic blok \mc{content}.
\end{itemize}

Jesli jednak funkcja zaczyna przejmowac odpowiedzialnosc za pol projektu naraz, moze to
byc znak, ze lepiej wydzielic osobny chunk albo osobny plik.

\section{Ograniczaj magię w wyrazeniach}
Bardzo długie warunki i skomplikowane wyrazenia wstawione bezposrednio do HTML-a wygladaja
sprytnie tylko przez chwile. Potem staja sie meczace.

Zamiast tego:

\begin{lstlisting}[style=moonchunk,caption={Lepsza czytelnosc przez nazwy pomocnicze}]
let canShowBanner: bool = isPublished and not isArchived;
let bannerLabel: string = canShowBanner ? "Widoczny" : "Ukryty";
\end{lstlisting}

\section{Utrzymuj pliki male i sensowne}
Jesli plik robi zbyt wiele rzeczy naraz, podziel go. Dobry podzial nie wynika z sztywnej
liczby linii, ale z odpowiedzialnosci. Przykładowo osobny plik moze miec:

\begin{itemize}
  \item wspolne metadata,
  \item logike bloga,
  \item zestaw chunkow sekcyjnych,
  \item dane zewnętrzne.
\end{itemize}

\section{Uzywaj \mc{print(...)} świadomie}
Debugowanie jest normalna częscią pracy, ale logi tez potrafia zasmiecic projekt.
Zostawiaj te, ktore rzeczywiscie pomagaja, i usuwaj te jednorazowe, gdy przestaja byc
potrzebne.

\section{Nie naduzywaj globali}
Globalne wartosci sa wygodne, ale jak kazde wygodne narzedzie mogą byc naduzywane.
Nie wszystko musi byc globalne. Trzymaj globalnie tylko to, co naprawde wspolne dla
wielu miejsc programu.

\section{Wybieraj prostszy typ, jesli wystarczy}
Nie trzeba wszedzie wpisywac \mc{unknown}, \mc{any} czy rozbudowanych rzutowan. Jesli
wystarcza \mc{string}, \mc{int}, \mc{array} albo \mc{object}, wybierz wlasnie je.
Kod bedzie prostszy i bezpieczniejszy.

\section{Zanim rozbudujesz, zbuduj maly pionowy przekroj}
Praktyczna rada projektowa: zanim napiszesz cala architekture, zbuduj maly kawalek,
ktory przechodzi przez wszystkie warstwy. Na przykład:

\begin{enumerate}
  \item jedna strona,
  \item jeden chunk wspolny,
  \item jedno zrodlo JSON,
  \item jedna petla,
  \item jeden warunek,
  \item jedna funkcja pomocnicza.
\end{enumerate}

Taki przekroj daje pewnosc, ze fundament dziala.

\section{Komentuj tam, gdzie jest decyzja, nie tam, gdzie jest oczywistosc}
Jesli dodajesz komentarz, niech wyjasnia \emph{dlaczego}, a nie tylko \emph{co}.

Dobry komentarz:

\begin{lstlisting}[style=moonchunk,caption={Komentarz wyjasniajacy decyzje}]
// Pokazujemy tylko wpisy opublikowane, bo szkice nie powinny trafiać do HTML-a.
if (posts[i].published) {
  ...
};
\end{lstlisting}

Slabszy komentarz to taki, ktory tylko powtarza kod.

\section{Buduj projekt tak, jakby mial urosnac}
Nawet jesli dzis projekt ma dwie podstrony, jutro moze miec dwadziescia. Nie chodzi o
przedwczesne komplikowanie, tylko o zdrowe podstawy:

\begin{itemize}
  \item sensowny podzial plikow,
  \item dobre nazwy,
  \item wspolne metadata,
  \item funkcje dla powtarzalnej logiki,
  \item dane ladowane z zewnatrz tam, gdzie to ma sens.
\end{itemize}

To zwykle wystarczy, aby projekt rozwijal sie spokojnie i bez chaosu.
